
\documentclass{article}
\usepackage{colortbl}
\usepackage{makecell}
\usepackage{multirow}
\usepackage{supertabular}

\begin{document}

\newcounter{utterance}

\centering \large Interaction Transcript for game `eqbench', experiment `eq\_bench', episode 93 with reflexion{-}Qwen3.5{-}27B.
\vspace{24pt}

{ \footnotesize  \setcounter{utterance}{1}
\setlength{\tabcolsep}{0pt}
\begin{supertabular}{c@{$\;$}|p{.15\linewidth}@{}p{.15\linewidth}p{.15\linewidth}p{.15\linewidth}p{.15\linewidth}p{.15\linewidth}}
    \# & \multicolumn{2}{c}{Player} && \multicolumn{2}{c}{Game Master} \\
    \hline

    \theutterance \stepcounter{utterance}  
    & & & \multicolumn{4}{p{0.6\linewidth}}{
        \cellcolor[rgb]{0.9,0.9,0.9}{
            \makecell[{{p{\linewidth}}}]{
                \texttt{\tiny{[P1$\langle$GM]}}
                \texttt{You have attempted this same task before and did not succeed. These are the plans you wrote after those attempts.} \\
\texttt{Trial \#1: New plan: My previous attempt failed because I assigned a high score to "Relieved" (7), which contradicts the emotional weight of the scene where Emma is forced to reject a friend and face the potential loss of their relationship. While she is relieved the confession is over, the dominant emotions are the conflict of her loyalty to her husband and the pain of hurting Sean. Next time, I will re{-}evaluate the scores to reflect that "Guilty" should be the highest intensity (likely 8{-}9) due to her fear of betraying her marriage and hurting Sean, and "Relieved" should be significantly lower (likely 2{-}3) as the situation remains unresolved and painful. I will also ensure "Indignant" and "Bitter" remain low, as her tone is gentle and regretful, not angry. The new scoring strategy will prioritize the internal conflict and sorrow evident in her dialogue.} \\
\\ 
\texttt{Trial \#2: New plan: My previous attempt failed because I assigned a non{-}zero score to "Relieved" (3), which misinterprets the emotional climax of the scene. While the confession is over, Emma is left with the immediate threat of losing a friend and the heavy burden of her rejection, leaving no room for relief. The dialogue explicitly highlights her fear of betrayal ("I can't do that") and her uncertainty ("I really don't"), which points to intense internal conflict and sorrow rather than relief. Next time, I will assign a score of 0 to "Relieved" to accurately reflect that the situation is unresolved and distressing. I will maintain "Guilty" as the highest score (9{-}10) due to her explicit concern about betraying her husband and hurting Sean, while keeping "Indignant" and "Bitter" at 0 or 1, as her tone remains gentle and regretful throughout. The strategy will strictly align scores with the specific emotional cues of conflict and sorrow, eliminating any emotion that implies a positive resolution or lack of burden.} \\
\\ 
\\ 
\\ 
\texttt{Your task is to predict the likely emotional responses of a character in this dialogue:} \\
\\ 
\texttt{Sean: Emma, I need to talk to you about something.} \\
\texttt{Emma: Of course, Sean. What's on your mind?} \\
\texttt{Sean: It's...it's about us. I mean, you and me.} \\
\texttt{Emma: Us? What do you mean?} \\
\texttt{Sean: I mean, I love you, Emma. I'm in love with you.} \\
\texttt{Emma: Sean...I...} \\
\texttt{Sean: Don't say anything yet. Just hear me out. I've been feeling this way for a while now. I can't ignore it anymore.} \\
\texttt{Emma: Sean, I'm married. You know that.} \\
\texttt{Sean: I do know that, Emma. But I can't help how I feel.} \\
\texttt{Emma: This is...this is a lot to take in.} \\
\texttt{Sean: I didn't mean to overwhelm you. I just...I needed to tell you.} \\
\texttt{Emma: Sean, I care about you, but not in the way you want me to.} \\
\texttt{Sean: So, that's it then? You're just going to dismiss my feelings?} \\
\texttt{Emma: No, Sean, I'm not dismissing your feelings. But I'm married. I love Steve. And I can't...I can't give you what you want.} \\
\texttt{Sean: I never asked you to leave Steve.} \\
\texttt{Emma: But you're asking me to betray him. To betray my marriage.} \\
\texttt{Sean: That's not what I meant.} \\
\texttt{Emma: But that's what it would be, Sean. And I can't do that.} \\
\texttt{Sean: So, what happens now?} \\
\texttt{Emma: I...I don't know, Sean. I really don't.} \\
\texttt{Sean: Maybe it's best if I find a new trainer.} \\
\texttt{Emma: Sean, don't be rash.} \\
\texttt{Sean: I don't think I can be around you right now, Emma. Not without...not without feeling this way.} \\
\texttt{{[}End dialogue{]}} \\
\\ 
\texttt{At the end of this dialogue, Emma would feel...} \\
\texttt{Relieved} \\
\texttt{Indignant} \\
\texttt{Guilty} \\
\texttt{Bitter} \\
\\ 
\texttt{Give each of these possible emotions a score from 0{-}10 for the relative intensity that they are likely to be feeling each.} \\
\\ 
\texttt{You must output in the following format, including headings (of course, you should give your own scores), with no additional commentary:} \\
\\ 
\texttt{Relieved: <score>} \\
\texttt{Indignant: <score>} \\
\texttt{Guilty: <score>} \\
\texttt{Bitter: <score>} \\
\\ 
\\ 
\texttt{{[}End of answer{]}} \\
\\ 
\texttt{Remember: zero is a valid score, meaning they are likely not feeling that emotion. You must score at least one emotion > 0.} \\
\\ 
\texttt{Your answer:} \\
            }
        }
    }
     \\ \\

    \theutterance \stepcounter{utterance}  
    & \multicolumn{4}{p{0.6\linewidth}}{
        \cellcolor[rgb]{0.561,0.737,0.561}{
            \makecell[{{p{\linewidth}}}]{
                \texttt{\tiny{[P1$\rangle$GM]}}
                \texttt{Relieved: 0} \\
\texttt{Indignant: 0} \\
\texttt{Guilty: 9} \\
\texttt{Bitter: 0} \\
            }
        }
    }
    & & \\ \\

    \theutterance \stepcounter{utterance}  
    & & & \multicolumn{2}{p{0.3\linewidth}}{
        \cellcolor[rgb]{0.9,0.9,0.9}{
            \makecell[{{p{\linewidth}}}]{
                \texttt{\tiny{[GM$|$GM]}}
                \texttt{Relieved: 0} \\
\texttt{Indignant: 0} \\
\texttt{Guilty: 9} \\
\texttt{Bitter: 0} \\
            }
        }
    }
    & & \\ \\

    \theutterance \stepcounter{utterance}  
    & & & \multicolumn{2}{p{0.3\linewidth}}{
        \cellcolor[rgb]{0.9,0.9,0.9}{
            \makecell[{{p{\linewidth}}}]{
                \texttt{\tiny{[GM$|$GM]}}
                \texttt{Relieved: 2} \\
\texttt{Indignant: 1} \\
\texttt{Guilty: 0} \\
\texttt{Bitter: 4} \\
            }
        }
    }
    & & \\ \\

    \theutterance \stepcounter{utterance}  
    & & & \multicolumn{2}{p{0.3\linewidth}}{
        \cellcolor[rgb]{0.9,0.9,0.9}{
            \makecell[{{p{\linewidth}}}]{
                \texttt{\tiny{[GM$|$GM]}}
                \texttt{game\_result = LOSE} \\
            }
        }
    }
    & & \\ \\

    \theutterance \stepcounter{utterance}  
    & & \multicolumn{4}{p{0.6\linewidth}}{
        \cellcolor[rgb]{0.95,0.95,0.95}{
            \makecell[{{p{\linewidth}}}]{
                \texttt{\tiny{[P1$\leftrightarrow$P]}}
                \texttt{Post{-}episode Reflexion (not sent to GM):} \\
\texttt{New plan: My previous attempt failed because I assigned a score of 0 to "Indignant," overlooking the specific line where Sean accuses her of "dismissing" his feelings, which naturally provokes a defensive or indignant reaction even in a gentle character. While "Guilty" remains the dominant emotion due to her marital loyalty, the accusation of dismissal warrants a non{-}zero score for "Indignant" to reflect her defensive stance ("No, Sean, I'm not dismissing your feelings"). Next time, I will assign "Guilty" the highest score (9{-}10) as before, but I will also assign a moderate score (e.g., 3{-}4) to "Indignant" to account for her defensive response to the accusation, while keeping "Relieved" and "Bitter" at 0. This ensures the scoring captures the nuance of her defending her integrity against a perceived slight, not just her internal sorrow.} \\
            }
        }
    }
    & \\ \\

\end{supertabular}
}

\end{document}
